\section{E1 - Phân loại ảnh: Softmax, MLP và CNN}

\subsection{Mục tiêu, dữ liệu và split}
E1 sử dụng Fashion-MNIST và giữ cùng dataset/split cho E2--E3 theo yêu cầu đề. Tập train gốc sẽ được tách train/validation bằng seed cố định; test set chỉ dùng cho đánh giá cuối.

\TODO{ĐIỀN SỐ MẪU TRAIN/VAL/TEST, TỶ LỆ TÁCH VÀ SEED SAU KHI CHẠY NOTEBOOK THẬT.}
\safeimage{e1/e1_eda.png}{EDA Fashion-MNIST: sample grid và phân bố 10 lớp.}

\subsection{Tiền xử lý và DataLoader}
Ảnh giữ cấu trúc grayscale $1\times28\times28$. Softmax/MLP flatten ảnh trong forward; CNN nhận tensor 2D trực tiếp. Báo cáo cần ghi batch size, normalization nếu dùng, số worker và lý do augmentation có/không được áp dụng.

\subsection{Kiến trúc bốn cấu hình}
\begin{table}[H]\centering
\caption{Các cấu hình E1.}
\begin{tabularx}{\textwidth}{@{}lllX@{}}
\toprule
ID & Mô hình & Kiến trúc chính & Câu hỏi so sánh \\
\midrule
E1-M1 & Linear Softmax & Flatten $\rightarrow$ Linear(784,10) & Baseline tuyến tính \\
E1-M2 & MLP & 784 $\rightarrow$ 256 $\rightarrow$ 128 $\rightarrow$ 10 & Giá trị của phi tuyến/depth \\
E1-M3 & CNN-2 & 2 Conv blocks $\rightarrow$ classifier & Inductive bias không gian \\
E1-M4 & CNN-3-BN & 3 Conv blocks + BatchNorm + Dropout & Capacity và regularization \\
\bottomrule
\end{tabularx}
\end{table}

\subsection{Training loop và protocol}
Training loop phải thể hiện tường minh \code{optimizer.zero\_grad()}, forward, criterion, \code{loss.backward()} và \code{optimizer.step()}. Tất cả mô hình dùng cùng split; hyperparameter khác biệt phải được ghi rõ thay vì ngầm thay đổi.

\begin{table}[H]\centering
\caption{Thiết lập huấn luyện E1.}
\begin{tabularx}{0.9\textwidth}{@{}lX@{}}
\toprule
Thành phần & Giá trị thực nghiệm \\
\midrule
Optimizer / LR & \TBD \\
Batch size / epochs & \TBD \\
Criterion & Cross-entropy \\
Seed & \TBD \\
Model selection & \TODO{CHỌN THEO VALIDATION ACCURACY/LOSS VÀ GHI RÕ} \\
\bottomrule
\end{tabularx}
\end{table}

\subsection{Kết quả định lượng}
\begin{table}[H]\centering
\caption{Kết quả E1 trên validation/test.}
\begin{tabular}{@{}lrrrrr@{}}
\toprule
Model & Params & Val Acc & Test Acc & Best epoch & Train time \\
\midrule
E1-M1 & \metricpending & \metricpending & \metricpending & \metricpending & \metricpending \\
E1-M2 & \metricpending & \metricpending & \metricpending & \metricpending & \metricpending \\
E1-M3 & \metricpending & \metricpending & \metricpending & \metricpending & \metricpending \\
E1-M4 & \metricpending & \metricpending & \metricpending & \metricpending & \metricpending \\
\bottomrule
\end{tabular}
\end{table}
\safeimage{e1/e1_learning_curves.png}{Learning curves của bốn cấu hình E1.}
\safeimage{e1/e1_confusion_errors.png}{Confusion matrix và các mẫu dự đoán sai tiêu biểu của E1.}

\subsection{Phân tích lỗi và kết luận E1}
\TODO{PHÂN TÍCH BẰNG EVIDENCE: lớp nào dễ nhầm, CNN cải thiện gì so với flatten, capacity có dẫn tới overfit không, tốc độ hội tụ và số tham số khác nhau thế nào.}
